\begin{problem}{1.09}
A hinge linkage is made of four thin rods of length $2l$ and four thin rods of length $l$; all the crossing points are connected by hinges, as shown in the figure. The hinge $O_3$ moves uniformly in the $x$ direction with speed $v$. When each rod makes an angle $\theta$ with the $x$ axis, find:
\begin{enumerate}[label=(\arabic*),leftmargin=2.2em,itemsep=1pt,topsep=2pt]
\item the magnitudes of the velocity and acceleration of the hinges $A_1$ and $B_3$;
\item the angular velocity and angular acceleration of the rod $A_3B_2$;
\item the equations of the trajectories of all the hinges.
\end{enumerate}
\end{problem}

\begin{figbox}{1.09}
\begin{tikzpicture}[scale=1.9]
  \def\c{0.643}
  \def\s{0.766}
  \coordinate (O)  at (0,0);
  \coordinate (A1) at (\c,\s);
  \coordinate (B1) at (\c,-\s);
  \coordinate (O1) at (2*\c,0);
  \coordinate (A2) at (3*\c,\s);
  \coordinate (B2) at (3*\c,-\s);
  \coordinate (O2) at (4*\c,0);
  \coordinate (A3) at (5*\c,\s);
  \coordinate (B3) at (5*\c,-\s);
  \coordinate (O3) at (6*\c,0);
  % axes
  \draw[axis] (-0.30,0) -- (6*\c+0.62,0) node[right] {$x$};
  \draw[axis] (0,-1.35) -- (0,1.35) node[above] {$y$};
  \foreach \yy in {-1.15,-0.85,...,1.15}
    \draw[thinline] (0,\yy) -- (-0.085,\yy+0.055);
  % rods
  \draw[thickline] (O)  -- (A1);
  \draw[thickline] (O)  -- (B1);
  \draw[thickline] (A1) -- (B2);
  \draw[thickline] (B1) -- (A2);
  \draw[thickline] (A2) -- (B3);
  \draw[thickline] (B2) -- (A3);
  \draw[thickline] (A3) -- (O3);
  \draw[thickline] (B3) -- (O3);
  % hinges
  \foreach \p in {O,A1,B1,O1,A2,B2,O2,A3,B3,O3}
    \draw[fill=white,line width=0.5pt] (\p) circle (1.7pt);
  % angle
  \draw[thinline] (0.34,0) arc (0:50:0.34);
  \node at (0.50,0.16) {$\theta$};
  % velocity of O3
  \draw[force] ($(O3)+(0,0.14)$) -- ++(0.62,0) node[above] {$v$};
  % labels
  \node[below left]  at (O)  {$O$};
  \node[below=2pt]   at (O1) {$O_1$};
  \node[below=2pt]   at (O2) {$O_2$};
  \node[below right] at (O3) {$O_3$};
  \node[above right] at (A1) {$A_1$};
  \node[above right] at (A2) {$A_2$};
  \node[above right] at (A3) {$A_3$};
  \node[below right] at (B1) {$B_1$};
  \node[below right] at (B2) {$B_2$};
  \node[below right] at (B3) {$B_3$};
\end{tikzpicture}
\end{figbox}

\begin{problem}{1.10}
A disc $A$ of radius $R$ rotates uniformly with angular velocity $\omega_1$ about its fixed axis $O_1$. Another disc $B$ of radius $r$ rolls without slipping along the rim of $A$, turning uniformly about its own axis $O_2$ with angular velocity $\omega_2$. Find the time required for $B$ to make one complete circuit around $A$.
\end{problem}

\begin{figbox}{1.10}
\begin{tikzpicture}[scale=1.75]
  \coordinate (O1) at (0,0);
  \coordinate (C)  at (40:1.9);
  \coordinate (O2) at (40:2.85);
  \draw[thickline] (O1) circle (1.9);
  \draw[thickline] (O2) circle (0.95);
  \draw[thinline] (O1) -- (O2);
  \fill (O1) circle (1.3pt);
  \fill (O2) circle (1.3pt);
  \node[below left] at (O1) {$O_1$};
  \node[right] at (O2) {$O_2$};
  \node at (-0.95,0.42) {$A$};
  \node at (2.55,1.45) {$B$};
  \node[above left] at (40:1.05) {$R$};
  \node[below,inner sep=1pt] at (2.10,1.50) {$r$};
  % sense of rotation
  \draw[force] (120:0.52) arc (120:340:0.52);
  \draw[force] (150:0.52) arc (150:-70:0.52);
\end{tikzpicture}
\end{figbox}

\begin{problem}{1.11}
A large steel ring of radius $R$ rotates in its own plane, in the counterclockwise sense, with constant angular velocity $\omega$ about a point $P$ on the ring. A small steel ring of radius $\frac13 R$ rolls without slipping along the inner side of the large ring in the same plane, in the direction shown in the figure. When the large ring turns once about $P$, the small ring turns exactly twice relative to the large ring. When the small ring has moved relative to the large ring to the position shown, find the accelerations of the points $A$ and $B$ of the small ring.
\end{problem}

\begin{figbox}{1.11}
\begin{tikzpicture}[scale=1.55]
  \def\R{2.0}
  \coordinate (O) at (0,0);
  \coordinate (P) at (-\R,0);
  \coordinate (A) at (0,-\R);
  \coordinate (B) at (0,-\R/3);
  \coordinate (S) at (0,-2*\R/3);
  \draw[thickline] (O) circle (\R);
  \draw[thickline] (S) circle (\R/3);
  \draw[dashedline] (P) -- (O) node[midway,above] {$R$};
  \draw[dashedline] (O) -- (A);
  \fill (O) circle (1.3pt);
  \fill (S) circle (1.3pt);
  % rotation of the large ring about P
  \draw[force] (P) ++(55:0.62) arc (55:160:0.62);
  \node at ($(P)+(105:0.85)$) {$\omega$};
  % radius of the small ring
  \draw[thinline] (S) -- ++(38:0.6667);
  \node at ($(S)+(0.30,-0.30)$) {$\frac{R}{3}$};
  % rolling sense
  \draw[force] (S) ++(55:0.44) arc (55:-75:0.44);
  \node[right] at (O) {$O$};
  \node[left] at (P) {$P$};
  \node[below] at (A) {$A$};
  \node[above right,inner sep=1pt] at (B) {$B$};
\end{tikzpicture}
\end{figbox}

\begin{problem}{1.12}
Between the walls of two coaxial thin-walled cylinders $A$ and $B$ of radii $R$ and $2R$, a small cylinder of radius exactly $\frac12 R$ is held. When $A$ and $B$ rotate uniformly in opposite directions with angular velocities $\omega_1$ and $\omega_2$, the small cylinder turns with them, there being no relative slipping at its points of contact with $A$ and $B$, as shown in the figure. Find:
\begin{enumerate}[label=(\arabic*),leftmargin=2.2em,itemsep=1pt,topsep=2pt]
\item the times required for the small cylinder to make one revolution relative to the ground and relative to cylinder $B$;
\item the acceleration of the point $C$ of the small cylinder that is in contact with $B$, at this instant, relative to the ground and relative to cylinder $A$.
\end{enumerate}
\end{problem}

\begin{figbox}{1.12}
\begin{tikzpicture}[scale=1.12]
  \def\R{1.55}
  \coordinate (O) at (0,0);
  \coordinate (S) at (0,-1.5*\R);
  \draw[thickline] (O) circle (\R);
  \draw[thickline] (O) circle (2*\R);
  \draw[thickline] (S) circle (0.5*\R);
  \draw[thinline] (O) -- (32:\R);
  \draw[thinline] (O) -- (-30:2*\R);
  \node[above right] at (32:0.55*\R) {$R$};
  \node[fill=white,inner sep=1pt] at (-30:1.42*\R) {$2R$};
  \fill (O) circle (1.3pt);
  % sense of rotation of A
  \draw[force] (238:0.60*\R) arc (238:300:0.60*\R);
  \node at (0,-0.33*\R) {$\omega_1$};
  % sense of rotation of B
  \draw[force] (0.55*\R,-2.05*\R) -- (-0.55*\R,-2.05*\R);
  \node at (0,-2.25*\R) {$\omega_2$};
  \node at (-0.52,0.55) {$A$};
  \node at (-2.35,0) {$B$};
  \node at (0,-1.83*\R) {$C$};
\end{tikzpicture}
\end{figbox}

\begin{problem}{1.13}
Two particles $A$ and $B$ move uniformly in the same sense along concentric circles of radii $a$ and $b$, their speeds being inversely proportional to the radii. Find the angle $\theta$ between the lines joining the two particles to the centre at the instant when the direction of their relative velocity is parallel to the line joining them.
\end{problem}

\begin{figbox}{1.13}
\begin{tikzpicture}[scale=1.6]
  \def\a{1.35}
  \def\b{2.05}
  \coordinate (O) at (0,0);
  \coordinate (A) at (24:\a);
  \coordinate (B) at (\b,0);
  \draw[thickline] (O) circle (\a);
  \draw[thickline] (O) circle (\b);
  \fill (A) circle (1.5pt);
  \fill (B) circle (1.5pt);
  \draw[thinline] (O) -- (A);
  \draw[thinline] (O) -- (B);
  \draw[thinline] (0.52,0) arc (0:24:0.52);
  \node at (0.68,0.15) {$\theta$};
  \node[above right] at (23:\a) {$A$};
  \node[right] at ($(B)+(0.10,0)$) {$B$};
  \node[above] at (0.68,0.42) {$a$};
  \node[below] at (1.62,0) {$b$};
  \node[below left] at (O) {$O$};
\end{tikzpicture}
\end{figbox}

\begin{problem}{1.14}
A smooth inclined rod $AB$ of length $l$ makes a fixed angle $\alpha$ with the vertical axis through its end $A$, and rotates about this axis with constant angular velocity $\omega$. A small ring of mass $m$ is threaded on the rod and starts from $A$ with an initial velocity $v_0$ relative to the rod, directed upward along the rod. Find:
\begin{enumerate}[label=(\arabic*),leftmargin=2.2em,itemsep=1pt,topsep=2pt]
\item the minimum value of $v_0$ for the ring to reach the end $B$;
\item the force exerted by the ring on the rod when the ring reaches $B$.
\end{enumerate}
\end{problem}

\begin{figbox}{1.14}
\begin{tikzpicture}[scale=1.35]
  \def\al{24}
  \def\L{3.5}
  \coordinate (A) at (0,0.45);
  \coordinate (B) at ({0.45+\L*sin(\al)},{0.45+\L*cos(\al)});
  \draw[thickline] (0,-0.15) -- (0,4.35);
  \draw[thickline] (A) -- (B);
  % axis rotation
  \draw[force] (0,4.12) ++(150:0.30) arc (150:-150:0.30);
  \node[right] at (0.28,4.12) {$\omega$};
  % angle alpha
  \draw[thinline] (0,1.10) arc (90:{90-\al}:0.65);
  \node at (0.20,1.28) {$\alpha$};
  \node[below left] at (A) {$A$};
  \node[above right] at (B) {$B$};
  % ring
  \coordinate (M) at ({0.45+0.60*\L*sin(\al)},{0.45+0.60*\L*cos(\al)});
  \draw[thickline] (M) circle (0.10);
  \draw[thickline,fill=white] (M) circle (0.055);
  \node[left] at ($(M)+(-0.20,-0.02)$) {$m$};
  % length l (dimension line parallel to the rod)
  \draw[thinline] (0.229,0.348) -- (2.103,3.546);
  \draw[thinline] (0.201,0.284) -- (0.257,0.412);
  \draw[thinline] (2.075,3.482) -- (2.131,3.610);
  \node[right] at (1.30,1.95) {$l$};
\end{tikzpicture}
\end{figbox}

\begin{problem}{1.15}
A smooth thin rod is bent into the shape shown in the figure and fixed in the vertical plane $(xOy)$, the equation of the curve being
\[
y=\dfrac{1}{2a}x^2
\]
where $a$ is a positive constant. A small ring of mass $m$ is threaded on the rod.
\begin{enumerate}[label=(\arabic*),leftmargin=2.2em,itemsep=1pt,topsep=2pt]
\item If the ring is released from rest at $x=2a$, find the force it exerts on the rod when it reaches $x=a$.
\item If the rod rotates about the $y$ axis with constant angular velocity $\omega_0$, for what value of $\omega_0$ can the ring remain at rest relative to the rod at any position on the rod?
\item If the angular velocity of the rod about the $y$ axis is taken to be $\frac{\sqrt2}{2}$ times the value found in (2), that is $\omega=\frac{\sqrt2}{2}\omega_0$, and the ring is released from rest relative to the rod at $x=2a$, what is the force it exerts on the rod when it reaches $x=a$?
\end{enumerate}
\end{problem}

\begin{figbox}{1.15}
\begin{tikzpicture}[scale=1.45]
  \draw[axis] (-2.55,0) -- (2.75,0) node[right] {$x$};
  \draw[axis] (0,-0.30) -- (0,3.10) node[above] {$y$};
  \draw[thickline,domain=-2.40:2.40,samples=100] plot (\x,{0.5*\x*\x});
  % marked point
  \coordinate (Q) at (2,2);
  \draw[thickline,fill=white] (Q) circle (0.10);
  \draw[dashedline] (Q) -- (2,0);
  \draw[thinline] (1,0) -- (1,0.10);
  \draw[thinline] (2,0) -- (2,0.10);
  \node[below] at (1,-0.03) {$a$};
  \node[below] at (2,-0.03) {$2a$};
  \node[below left] at (0,0) {$O$};
\end{tikzpicture}
\end{figbox}

\begin{problem}{1.16}
A large horizontal disc of radius $R$ rotates uniformly with angular velocity $\Omega$ about the vertical axis through its centre. A small disc of radius $r$ and mass $m$ is gently placed on the large disc from above. After the small disc has settled, the distance between the axes of the two discs is $d$ ($d>r$, $R>d+r$). The coefficient of friction between the discs is $\mu$.
\begin{enumerate}[label=(\arabic*),leftmargin=2.2em,itemsep=1pt,topsep=2pt]
\item Find the angular velocity $\omega$ of the small disc in its steady rotation.
\item What torque must be applied to the large disc to keep its angular speed unchanged?
\end{enumerate}
\end{problem}

\begin{figbox}{1.16}
\begin{tikzpicture}[scale=1.15]
  \def\Rx{3.2}
  \def\Ry{1.15}
  \draw[thickline] (0,0) ellipse ({\Rx} and {\Ry});
  \draw[thickline] (0,-0.10) ellipse ({\Rx} and {\Ry});
  \draw[thickline] (0,0) -- (0,2.95);
  \draw[force] (0,2.95) ++(150:0.35) arc (150:-150:0.35);
  \node[right] at (0.30,2.95) {$\Omega$};
  % radius R
  \draw[thinline] (0,0) -- (-2.62,-0.66);
  \node[above right] at (-1.32,-0.32) {$R$};
  % small disc
  \coordinate (S) at (1.75,0.30);
  \draw[thickline] (S) ellipse (0.55 and 0.27);
  \draw[thickline] (S) -- ++(0,0.42);
  \node[right] at ($(S)+(0.28,0.12)$) {$r$};
  \draw[dim] (0,0.74) -- (1.75,0.74);
  \node[above] at (0.87,0.78) {$d$};
\end{tikzpicture}
\end{figbox}

\begin{problem}{1.17}
Two identical small balls $A$ and $B$, each of mass $M$ and radius $R$, are suspended from the same point $O$ by two strings each of length $l=2R$. A small ball $C$ of mass $m(=nM)$ and radius $r\left(=\frac12R\right)$ is placed on the two balls. The surfaces of all three balls are smooth. Discuss the relation between the angle $\theta$ that the strings make with the vertical and $n$ when the system is in equilibrium.
\end{problem}

\begin{figbox}{1.17}
\begin{tikzpicture}[scale=1.5]
  \coordinate (O)  at (0,4.30);
  \coordinate (CA) at (-1.15,0.95);
  \coordinate (CB) at (1.15,0.95);
  \coordinate (CC) at (0,2.236);
  % ceiling
  \draw[thickline] (-2.9,4.30) -- (2.9,4.30);
  \foreach \xx in {-2.8,-2.55,...,2.8}
    \draw[thinline] (\xx,4.30) -- ++(-0.10,0.17);
  % strings
  \draw[thickline] (O) -- (-0.766,2.034);
  \draw[thickline] (O) -- (0.766,2.034);
  % balls
  \draw[thickline] (CA) circle (1.15);
  \draw[thickline] (CB) circle (1.15);
  \draw[thickline] (CC) circle (0.575);
  % vertical through O and angle theta
  \draw[dashedline] (O) -- (0,1.25);
  \draw[thinline] ($(O)+(0,-0.95)$) arc (-90:-109.5:0.95);
  \node at (-0.20,3.40) {$\theta$};
  % radii
  \draw[thinline] (CA) -- ++(215:1.15);
  \node[above] at (-1.85,0.50) {$R$};
  \draw[thinline] (CB) -- ++(-25:1.15);
  \node[above] at (1.95,0.60) {$R$};
  \draw[thinline] (CC) -- ++(40:0.575);
  \node[fill=white,inner sep=1pt] at (0.43,2.55) {$r$};
  % labels
  \node at (-1.08,0.60) {$M$};
  \node at (1.32,0.55) {$M$};
  \node[left] at (-0.22,2.35) {$m$};
  \node[below] at (0,1.95) {$C$};
  \node[below] at (-1.15,-0.28) {$A$};
  \node[below] at (1.15,-0.28) {$B$};
  \node[left] at (-0.60,3.10) {$l$};
  \fill (O) circle (1.3pt);
  \node[above right] at (O) {$O$};
\end{tikzpicture}
\end{figbox}

\begin{problem}{1.18}
A small ball of mass $m$ is fixed to a large ring of mass $M$. The ring is hung on a rough nail, as shown in the figure. If the ring is to remain in equilibrium when hung on the nail at any point of the ring (except the point at which the small ball is fixed), what is the minimum coefficient of friction $\mu$ between the ring and the nail?
\end{problem}

\begin{figbox}{1.18}
\begin{tikzpicture}[scale=1.6]
  \draw[thickline] (0,0) circle (1.55);
  \coordinate (O) at (58:1.55);
  \coordinate (m) at (-25:1.55);
  \draw[thickline] (O) -- ++(58:0.32);
  \fill (O) circle (2.2pt);
  \node[left] at ($(O)+(-0.28,0.04)$) {$O$};
  \fill (m) circle (2.2pt);
  \node[right] at ($(m)+(0.12,0)$) {$m$};
  \node at (-0.32,-0.02) {$M$};
\end{tikzpicture}
\end{figbox}

\begin{problem}{1.19}
Two small balls 1 and 2 of equal mass, connected by a light string, are placed on an inclined plane of angle $\theta$. The coefficients of friction between the plane and balls 1 and 2 are $\mu_1$ and $\mu_2$ respectively, with $\mu_1<\mu_2$ and $\tan\theta=\sqrt{\mu_1\mu_2}$. Find the maximum value of the angle $\alpha$ between the string and the line of greatest slope of the plane for the system to be in equilibrium.
\end{problem}

\begin{figbox}{1.19}
\begin{tikzpicture}[scale=1.45]
  \coordinate (V)  at (0,0);
  \coordinate (Va) at (3.40,0.55);
  \coordinate (Vb) at (1.45,2.75);
  \draw[thickline] (V) -- (Va) -- ($(Va)+(Vb)$) -- (Vb) -- cycle;
  \coordinate (P1) at (1.678,1.177);
  \coordinate (P2) at (2.740,1.950);
  \draw[thickline] (P1) -- (P2);
  \draw[thickline,fill=white] (P1) circle (0.11);
  \draw[thickline,fill=white] (P2) circle (0.11);
  \draw[dashedline] (2.484,1.464) -- (2.996,2.436);
  \draw[thinline] (2.400,1.703) arc (216:242:0.42);
  \node at (2.28,1.44) {$\alpha$};
  \draw[thinline] (0.790,0.128) arc (9.2:62.2:0.80);
  \node at (0.95,0.46) {$\theta$};
  \node[above left] at (P1) {$1$};
  \node[above right] at (P2) {$2$};
\end{tikzpicture}
\end{figbox}

\begin{problem}{1.20}
A hinge structure is made of four uniform rods of length $2l$ and mass $2m$ and four uniform rods of length $l$ and mass $m$. The structure has ten light smooth hinges. Hinge $O_1$ is suspended from a horizontal axis, and a string of length $\sqrt2\,l$ joins hinges $O_3$ and $O_4$. When the structure is in equilibrium, the angle between the two rods at every hinge is $90\degree$, as shown in the figure. Find the tension in the string in equilibrium.
\end{problem}

\begin{figbox}{1.20}
\begin{tikzpicture}[scale=1.30]
  \coordinate (O1) at (0,4.80);
  \coordinate (A1) at (-0.85,3.95);
  \coordinate (B1) at (0.85,3.95);
  \coordinate (O2) at (0,3.10);
  \coordinate (A2) at (-0.85,2.25);
  \coordinate (B2) at (0.85,2.25);
  \coordinate (O3) at (0,1.40);
  \coordinate (A3) at (-0.85,0.55);
  \coordinate (B3) at (0.85,0.55);
  \coordinate (O4) at (0,-0.30);
  % ceiling
  \draw[thickline] (-1.8,4.95) -- (1.8,4.95);
  \foreach \xx in {-1.7,-1.45,...,1.7}
    \draw[thinline] (\xx,4.95) -- ++(-0.10,0.17);
  % rods
  \draw[thickline] (O1) -- (A1);
  \draw[thickline] (O1) -- (B1);
  \draw[thickline] (A1) -- (B2);
  \draw[thickline] (B1) -- (A2);
  \draw[thickline] (A2) -- (B3);
  \draw[thickline] (B2) -- (A3);
  \draw[thickline] (A3) -- (O4);
  \draw[thickline] (B3) -- (O4);
  % string
  \draw[thinline] (O3) -- (O4);
  % hinges
  \foreach \p in {O1,A1,B1,O2,A2,B2,O3,A3,B3,O4}
    \draw[fill=white,line width=0.5pt] (\p) circle (1.7pt);
  % labels
  \node[right] at ($(O1)+(0.10,-0.10)$) {$O_1$};
  \node[left]  at ($(A1)+(-0.10,0)$) {$A_1$};
  \node[right] at ($(B1)+(0.10,0)$) {$B_1$};
  \node[right] at ($(O2)+(0.10,-0.10)$) {$O_2$};
  \node[left]  at ($(A2)+(-0.10,0)$) {$A_2$};
  \node[right] at ($(B2)+(0.10,0)$) {$B_2$};
  \node[right] at ($(O3)+(0.10,-0.10)$) {$O_3$};
  \node[left]  at ($(A3)+(-0.10,0)$) {$A_3$};
  \node[right] at ($(B3)+(0.10,0)$) {$B_3$};
  \node[right] at ($(O4)+(0.10,-0.10)$) {$O_4$};
\end{tikzpicture}
\end{figbox}

\begin{problem}{1.21}
At the foot of an inclined plane of angle $\theta$, a vertical board of length $l$ is connected by a smooth hinge. A sphere of radius $R$ and mass $m$ is held between the board and the plane, and a horizontal force $F$ perpendicular to the board is applied at the upper end of the board so that both the sphere and the board remain in equilibrium, as shown in the figure. The coefficient of friction between the board and the sphere is $\mu_1$, and that between the sphere and the plane is $\mu_2$. The force $F$ is now gradually increased. At what value of $F$ is the equilibrium of the system destroyed?
\end{problem}

\begin{figbox}{1.21}
\begin{tikzpicture}[scale=0.95]
  % ground
  \draw[thickline] (-0.9,0) -- (6.4,0);
  \foreach \xx in {-0.8,-0.5,...,6.3}
    \draw[thinline] (\xx,0) -- ++(-0.12,-0.12);
  % incline
  \draw[thickline] (0,0) -- (6.0,3.464);
  % board
  \draw[thickline,fill=white] (-0.045,0) rectangle (0.045,3.80);
  \fill (0,0) circle (1.6pt);
  % ball
  \draw[thickline] (1.05,1.82) circle (1.05);
  \fill (1.05,1.82) circle (1.3pt);
  \draw[thinline] (1.05,1.82) -- ++(-40:1.05);
  \node[above right] at (1.26,1.58) {$R$};
  % angle theta
  \draw[thinline] (0.85,0) arc (0:30:0.85);
  \node at (1.05,0.24) {$\theta$};
  % force F
  \draw[force] (-1.30,3.60) -- (-0.06,3.60);
  \node[above] at (-0.72,3.66) {$F$};
  % length l
  \draw[dim] (-0.55,0.05) -- (-0.55,3.35);
  \node[left] at (-0.60,1.70) {$l$};
\end{tikzpicture}
\end{figbox}

\begin{problem}{1.22}
Three identical uniform rods, each of mass $m$ and length $l$, lean against one another as shown in the figure; the lines joining their points of contact with the ground form an equilateral triangle of side $l$. The coefficients of friction between the rods and the ground are equal.
\begin{enumerate}[label=(\arabic*),leftmargin=2.2em,itemsep=1pt,topsep=2pt]
\item Find the magnitude and direction of the force acting on the upper end of rod $OA$.
\item If a small ball of mass $m$ is also fixed at the midpoint of rod $OA$, find again the magnitude and direction of the force acting on its upper end.
\item Find the minimum coefficient of friction between the rods and the ground for the system to remain at rest.
\end{enumerate}
\end{problem}

\begin{figbox}{1.22}
\begin{tikzpicture}[scale=1.35]
  \coordinate (O) at (0,3.55);
  \coordinate (A) at (-1.90,0);
  \coordinate (B) at (1.90,0);
  \coordinate (C) at (0.35,1.10);
  \draw[dashedline] (A) -- (B) -- (C) -- cycle;
  \draw[thickline] (O) -- (A);
  \draw[thickline] (O) -- (B);
  \draw[thickline] (O) -- (C);
  \node[above] at (O) {$O$};
  \node[below left] at (A) {$A$};
  \node[below right] at (B) {$B$};
  \node[right] at (C) {$C$};
\end{tikzpicture}
\end{figbox}

\begin{problem}{1.23}
A ring of radius $R$ stands upright on the ground at point $P$. The line $AP$ is perpendicular to the plane of the ring, and $\overline{AP}=\frac32R$. One end of a thin rod is hinged at $A$, and the body of the rod rests on the ring at point $B$; the radius $OB$ through $B$ makes an angle $\theta=60\degree$ with $OP$, as shown in the figure. Find the minimum coefficient of friction between the rod and the ring for the rod to be in equilibrium in this position.
\end{problem}

\begin{figbox}{1.23}
\begin{tikzpicture}[scale=1.15,yscale=0.95]
  \coordinate (O) at (0,0);
  \coordinate (P) at (0,-2.20);
  \coordinate (B) at (1.342,-1.100);
  \coordinate (A) at (3.05,-3.35);
  \coordinate (Bp) at (1.008,-0.660);
  \draw[thickline] (O) ellipse (1.55 and 2.20);
  \draw[dashedline] (O) -- (P);
  \draw[dashedline] (O) -- (B);
  \node[left] at (-0.14,0.10) {$O$};
  \node[below left] at (P) {$P$};
  \node[above right] at (1.38,-1.08) {$B$};
  \node[right] at (A) {$A$};
  \node[left] at (-0.28,-1.25) {$R$};
  \draw[thinline] ($(O)+(-90:0.85)$) arc (-90:-39.3:0.85);
  \node at (0.30,-1.60) {$\theta$};
  % rod
  \draw[thickline] (Bp) -- (A);
  % AP
  \draw[thinline] (P) -- (A);
  % right angle at P
  \draw[thinline] (0,-2.02) -- (0.165,-1.972) -- (0.151,-2.190);
\end{tikzpicture}
\end{figbox}

\begin{problem}{1.24}
A jet of water from a fountain holds a rubbish bin of mass $M$ upside down in the air. The water leaves the fountain nozzle of area $S_0$ with speed $v_0$ and is directed vertically upward. After striking the inner bottom of the bin, half of the water is retained on the bottom and flows down along the inner wall, while the other half splashes vertically downward with its original speed, as shown in the figure. Find the height $h$ at which the bin remains in the air. The density of water is $\rho$.
\end{problem}

\begin{figbox}{1.24}
\begin{tikzpicture}[scale=0.95]
  % ground
  \draw[thickline] (-2.7,0) -- (3.4,0);
  \foreach \xx in {-2.6,-2.3,...,3.3}
    \draw[thinline] (\xx,0) -- ++(-0.12,-0.12);
  % walls of the inverted bin
  \draw[thickline] (-0.95,4.25) -- (-1.40,1.60);
  \draw[thickline] (0.95,4.25) -- (1.40,1.60);
  \draw[thickline] (-1.02,4.25) -- (1.02,4.25);
  \draw[thickline] (-1.02,4.15) -- (1.02,4.15);
  % jet
  \draw[thickline] (0,0.08) -- (0,4.13);
  \draw[force] (0,0.30) -- (0,0.90);
  \node[right] at (0.10,0.58) {$v_0$};
  % water spreading on the bottom and running down the walls
  \draw[thickline] (0,4.13) .. controls (0.26,4.02) and (0.68,4.00) .. (0.80,3.88)
        .. controls (1.00,3.55) and (1.10,2.60) .. (1.12,1.78);
  \draw[thickline] (0,4.13) .. controls (-0.26,4.02) and (-0.68,4.00) .. (-0.80,3.88)
        .. controls (-1.00,3.55) and (-1.10,2.60) .. (-1.12,1.78);
  % water splashing back downward
  \draw[dashedline] (0,4.10) -- (0,1.02);
  \draw[dashedline] (0.08,4.10) -- (0.08,1.20);
  % droplets
  \foreach \p in {(0.52,2.90),(-0.48,2.55),(0.40,2.10),(-0.55,3.20),(0.25,1.55),(-0.20,1.35)}
    \fill \p circle (0.030);
  % height h
  \draw[dim] (2.60,0) -- (2.60,4.25);
  \draw[thinline] (2.40,4.25) -- (2.95,4.25);
  \draw[thinline] (2.40,0) -- (2.95,0);
  \node[right] at (2.67,2.12) {$h$};
  \node[below] at (0,-0.22) {$S_0$};
\end{tikzpicture}
\end{figbox}

\begin{problem}{1.25}
A very small spherical raindrop condenses as a nucleus in a uniform, still cloud layer and falls from rest, capturing all the droplets it meets on its way, so that its volume increases continuously while it remains spherical. Neglecting the viscous drag on the drop, show that the drop tends to fall with uniform acceleration, and find this acceleration.
\end{problem}

\begin{problem}{1.26}
A block of mass $m$ is connected to a soft rope of linear mass density $\lambda$. Initially the block is at the top of an inclined plane of angle $\theta$, and the rope is coiled on the platform at the upper edge of the plane. The plane and the platform are both smooth. The block is released and slides down the plane, as shown in the figure. Find the speed of $m$ after it has slid a distance $x$ down the plane ($x$ less than the length of the rope).
\end{problem}

\begin{figbox}{1.26}
\begin{tikzpicture}[scale=0.92]
  % wedge
  \draw[thickline] (0,0) -- (5.50,2.00) -- (7.40,2.00) -- (7.40,0) -- cycle;
  \draw[thickline] (-0.3,0) -- (7.6,0);
  \foreach \xx in {-0.2,0.1,...,7.5}
    \draw[thinline] (\xx,0) -- ++(-0.12,-0.12);
  % angle theta
  \draw[thinline] (0.90,0) arc (0:20:0.90);
  \node at (1.12,0.20) {$\theta$};
  % block on the incline
  \begin{scope}[shift={(2.44,0.889)},rotate=20]
    \draw[thickline,fill=white] (-0.30,-0.02) rectangle (0.30,0.42);
    \node at (0,0.20) {$m$};
  \end{scope}
  % rope up the incline and coiled on the platform
  \draw[thickline] (2.58,1.39) -- (5.42,2.02) -- (6.05,2.02);
  \draw[thickline] (6.05,2.10) ellipse (0.18 and 0.075);
  \draw[thickline] (6.42,2.09) ellipse (0.18 and 0.075);
  \draw[thickline] (6.79,2.11) ellipse (0.18 and 0.075);
\end{tikzpicture}
\end{figbox}

\begin{problem}{1.27}
A small ball of mass $m$ lies on a table, and a vertical light spring of stiffness $k$ is attached to it. Initially the spring is at its natural length; its upper end is then made to move vertically upward with constant speed $v$ by an external force, as shown in the figure. Find the work done by the force acting on the upper end of the spring during the interval from the start of the motion of that end until the spring first reaches its maximum extension.
\end{problem}

\begin{figbox}{1.27}
\begin{tikzpicture}[scale=1.25]
  % table
  \draw[thickline] (-1.6,0) -- (1.6,0);
  \foreach \xx in {-1.5,-1.2,...,1.5}
    \draw[thinline] (\xx,0) -- ++(-0.12,-0.12);
  % ball
  \draw[thickline] (0,0.30) circle (0.30);
  \node[right] at (0.36,0.30) {$m$};
  % spring
  \draw[thickline,decorate,decoration={coil,aspect=0.5,segment length=4.5pt,amplitude=4.5pt}]
    (0,0.60) -- (0,2.55);
  \node[right] at (0.18,1.55) {$k$};
  % upper end moving up with speed v
  \draw[thickline] (0,2.55) -- (0,3.05);
  \draw[force] (0,3.05) -- (0,3.55);
  \node[right] at (0.08,3.35) {$v$};
\end{tikzpicture}
\end{figbox}
